\section{Objectius}

L'objectiu general del treball és \textbf{dissenyar i avaluar una metodologia industrialitzable de Root Cause Analysis basada en Machine Learning} aplicable a entorns farmacèutics regulats.

Per assolir aquest objectiu general, es defineixen els objectius específics següents:

\begin{itemize}
  \item \textbf{Objectiu 1.} Definir un marc conceptual de RCA data-driven alineat amb requisits GxP (traçabilitat, integritat de dades, validació i governança).
  \item \textbf{Objectiu 2.} Construir un pipeline de dades reproduïble per integrar variables de procés, qualitat i context operatiu rellevants per a l'anàlisi de desviacions.
  \item \textbf{Objectiu 3.} Desenvolupar i comparar models de ML interpretables i models de major capacitat per predir esdeveniments de desviació.
  \item \textbf{Objectiu 4.} Aplicar tècniques d'explicabilitat (global i local) per prioritzar hipòtesis de causa arrel amb criteris quantitatius.
  \item \textbf{Objectiu 5.} Proposar una arquitectura d'industrialització i un pla de validació conceptual que permetin la integració progressiva del sistema en operació.
\end{itemize}

Com a objectius complementaris, es planteja també: (i) quantificar l'impacte operatiu de la metodologia sobre l'eficiència de la investigació; (ii) identificar riscos ètics i regulatoris del model; i (iii) establir un conjunt de mètriques combinades que incloguin tant rendiment predictiu com utilitat per a RCA.
